\documentclass[a4paper,12pt]{article}

\usepackage{fontspec}
\usepackage[vietnamese,provide=*]{babel}
\usepackage{graphicx}
\usepackage{booktabs}
\usepackage{minted}
\usepackage[hidelinks]{hyperref}

\setmainfont{Latin Modern Roman}
\setsansfont{Latin Modern Sans}
\setmonofont{Latin Modern Mono}

\title{Báo cáo bài tập E1--E3}
\author{TODO: tên thành viên}
\date{HK261}

\begin{document}
\maketitle
\newpage
\tableofcontents
\newpage

\section*{Liêm chính học thuật}
\addcontentsline{toc}{section}{Liêm chính học thuật}
Nhóm cam kết tuân thủ quy định về liêm chính học thuật của nhà trường. Công cụ AI chỉ được sử dụng ở mức hỗ trợ. Toàn bộ mã nguồn, thí nghiệm và số liệu trong báo cáo do nhóm thực hiện. Các thành viên hiểu và có thể giải thích đầy đủ các nội dung đã nộp.

\section{Giới thiệu}\label{sec:intro}

\subsection{Phạm vi và mục tiêu}\label{sec:scope}

Báo cáo này trình bày các bài tập E1--E3 của môn học. Cả ba bài tập đều giải bài toán phân loại ảnh trên cùng một tập dữ liệu nhỏ (MNIST) với cùng cách chia dữ liệu, nhằm so sánh các họ mô hình ở điều kiện công bằng.

Mục tiêu của từng bài tập như sau:

\begin{itemize}
    \item \textbf{E1} (Mục~\ref{sec:e1}): xây dựng và so sánh ba bộ phân loại, gồm có Softmax, Multilayer Perceptron (MLP) và Convolutional Neural Networks (CNN).
    \item \textbf{E2} (Mục~\ref{sec:e2}): hiện thực Multi-Head Self-Attention (MSA) từ đầu, đối chiếu với bản \texttt{nn.MultiheadAttention} của PyTorch, và so sánh ít nhất hai cách tách ảnh thành token.
    \item \textbf{E3} (Mục~\ref{sec:e3}): mô hình hóa ảnh như dữ liệu chuỗi bằng Long Short-Term Memory (LSTM) và Gated Recurrent Unit (GRU), thử các cách biểu diễn chuỗi khác nhau và so sánh với baseline MLP và CNN của E1.
\end{itemize}

Toàn bộ mã nguồn, cấu hình huấn luyện được công bố trên GitHub của nhóm với hướng dẫn môi trường và lệnh chạy.
\begin{itemize}
    \item Repository: \url{https://github.com/tronglh241/co5085_excersises_and_assignments}
    \item GitHub Pages: \url{https://tronglh241.github.io/co5085_excersises_and_assignments}
\end{itemize}
Báo cáo này cũng được đặt trong repository tại \texttt{reports/exercises.pdf}.

\subsection{Ký hiệu}
% TODO
Bảng~\ref{tab:notation} liệt kê các ký hiệu dùng xuyên suốt báo cáo.

\begin{table}[ht]
    \centering
    \caption{Ký hiệu sử dụng trong báo cáo.}
    \label{tab:notation}
    \begin{tabular}{ll}
        \toprule
        Ký hiệu                                                          & Ý nghĩa                                                  \\
        \midrule
        $\mathbf{x} \in \mathbb{R}^{1 \times H \times W}$                & Ảnh đầu vào ($H = W = 28$)                               \\
        $y \in \{0, \dots, C-1\}$                                        & Nhãn lớp, $C = 10$                                       \\
        $N$, $B$                                                         & Số mẫu của tập dữ liệu, kích thước batch                 \\
        $\boldsymbol{\theta}$                                            & Tham số của mô hình                                      \\
        $f_{\boldsymbol{\theta}}(\mathbf{x}) \in \mathbb{R}^{C}$         & Logit đầu ra của mô hình                                 \\
        $\mathrm{softmax}(\mathbf{z})_k = e^{z_k} / \sum_j e^{z_j}$      & Hàm softmax                                              \\
        $\mathcal{L}$                                                    & Hàm mất mát (cross-entropy)                              \\
        \midrule
        $T$                                                              & Số token (hoặc số bước thời gian)                        \\
        $d$                                                              & Số chiều embedding của token                             \\
        $n_h$, $d_h = d / n_h$                                           & Số head và số chiều mỗi head                             \\
        $\mathbf{Q}, \mathbf{K}, \mathbf{V}$                             & Ma trận query, key, value                                \\
        $\mathbf{W}^{Q}, \mathbf{W}^{K}, \mathbf{W}^{V}, \mathbf{W}^{O}$ & Các ma trận chiếu của MSA                                \\
        \midrule
        $\mathbf{x}_t$                                                   & Đầu vào tại bước thời gian $t$ (một hàng hoặc một patch) \\
        $\mathbf{h}_t$, $\mathbf{c}_t$                                   & Trạng thái ẩn và trạng thái ô nhớ (LSTM) tại bước $t$    \\
        \bottomrule
    \end{tabular}
\end{table}

\section{Thiết lập thực nghiệm chung}\label{sec:setup}

Mục này mô tả các thiết lập dùng chung cho cả ba bài tập, bao gồm: tập dữ liệu, cách chia dữ liệu, tiền xử lý, quy trình huấn luyện, độ đo và cách tái lập thí nghiệm. Các mô hình ở E1--E3 đều được huấn luyện và đánh giá trên cùng một cách chia dữ liệu, do đó kết quả giữa các bài tập có thể so sánh trực tiếp.

\subsection{Tập dữ liệu}

Tập dữ liệu mà nhóm sử dụng là tập \textbf{MNIST}, chứa các hình ảnh chữ số viết tay từ 0 đến 9. Các hình ảnh là ảnh xám có kích thước $28 \times 28$ với giá trị điểm ảnh trong khoảng $[0, 255]$. Tập dữ liệu gốc gồm tổng cộng 70\,000 ảnh, trong đó, 60\,000 ảnh thuộc tập huấn luyện và 10\,000 ảnh thuộc tập kiểm tra. Đây là một tập nhỏ, nằm trong danh sách tập dữ liệu được đề bài cho phép, phù hợp để so sánh nhiều họ mô hình với chi phí huấn luyện thấp.

Hình~\ref{fig:mnist-samples} minh họa các mẫu của từng lớp được lấy ra từ tập huấn luyện của \textbf{MNIST}.

\begin{figure}[ht]
    \centering
    \includegraphics[width=0.8\linewidth]{images/e1_e3/samples}
    \caption{Một số mẫu ngẫu nhiên của từng lớp trong MNIST.}
    \label{fig:mnist-samples}
\end{figure}

\subsection{Chia tập dữ liệu}\label{sec:split}

Nhóm giữ nguyên tập kiểm tra chính thức của MNIST và chỉ dùng nó cho lần đánh giá cuối cùng. Tập validation được lấy ra từ 60\,000 ảnh huấn luyện gốc: 10\,000 ảnh dành cho validation, 50\,000
ảnh còn lại dùng để huấn luyện (Bảng~\ref{tab:split}). Validation giúp theo
dõi quá trình huấn luyện, điều chỉnh learning rate và chọn checkpoint tốt
nhất mà không chạm vào tập test.

Việc tách train/validation được thực hiện theo từng lớp (\textit{stratified}),
nên tỷ lệ các chữ số ở hai tập này gần như trùng nhau
(Hình~\ref{fig:class-dist}, Bảng~\ref{tab:class-dist}). Quá trình tách dùng
seed cố định (\texttt{\todo{seed}}), vì vậy mọi lần chạy ở E1, E2 và E3 đều
cho đúng cùng một cách chia. Nhờ đó các mô hình được so sánh trên cùng dữ
liệu.

Các lớp phân bố khá đồng đều. Chữ số 1 chiếm nhiều nhất (11{,}24\% tập train),
chữ số 5 ít nhất (9{,}04\%), chênh lệch khoảng 1{,}24 lần. Mức lệch này nhỏ,
nên nhóm dùng accuracy làm độ đo chính và không cần xử lý mất cân bằng lớp.

\begin{table}[ht]
    \centering
    \caption{Cách chia dữ liệu. Tỷ lệ tính trên tổng 70\,000 ảnh.}
    \label{tab:split}
    \begin{tabular}{lrr}
        \toprule
        Tập        & Số ảnh  & Tỷ lệ    \\
        \midrule
        Train      & 50\,000 & 71{,}4\% \\
        Validation & 10\,000 & 14{,}3\% \\
        Test       & 10\,000 & 14{,}3\% \\
        \midrule
        Tổng       & 70\,000 & 100\%    \\
        \bottomrule
    \end{tabular}
\end{table}

\begin{table}[ht]
    \centering
    \caption{Tỷ lệ (\%) của từng chữ số trong mỗi tập. Train và validation được
        tách phân tầng nên gần như trùng nhau; tập test là tập chính thức của MNIST.}
    \label{tab:class-dist}
    \small
    \begin{tabular}{lrrrrrrrrrr}
        \toprule
        Tập        & 0      & 1       & 2       & 3       & 4      & 5      & 6      & 7       & 8      & 9       \\
        \midrule
        Train      & 9{,}87 & 11{,}24 & 9{,}93  & 10{,}22 & 9{,}74 & 9{,}04 & 9{,}86 & 10{,}44 & 9{,}75 & 9{,}91  \\
        Validation & 9{,}87 & 11{,}24 & 9{,}93  & 10{,}22 & 9{,}74 & 9{,}03 & 9{,}86 & 10{,}44 & 9{,}75 & 9{,}92  \\
        Test       & 9{,}80 & 11{,}35 & 10{,}32 & 10{,}10 & 9{,}82 & 8{,}92 & 9{,}58 & 10{,}28 & 9{,}74 & 10{,}09 \\
        \bottomrule
    \end{tabular}
\end{table}

\begin{figure}[ht]
    \centering
    \begin{tikzpicture}
        \begin{axis}[
                ybar,
                width=\linewidth, height=6cm,
                bar width=4.5pt,
                xtick={0,...,9},
                xlabel={Chữ số}, ylabel={Tỷ lệ (\%)},
                ymin=0, ymax=13,
                enlarge x limits=0.08,
                ymajorgrids,
                legend style={at={(0.5,1.03)}, anchor=south, legend columns=-1},
            ]
            \addplot coordinates {(0,9.87) (1,11.24) (2,9.93)  (3,10.22) (4,9.74)
                    (5,9.04) (6,9.86)  (7,10.44) (8,9.75)  (9,9.91)};
            \addplot coordinates {(0,9.87) (1,11.24) (2,9.93)  (3,10.22) (4,9.74)
                    (5,9.03) (6,9.86)  (7,10.44) (8,9.75)  (9,9.92)};
            \addplot coordinates {(0,9.80) (1,11.35) (2,10.32) (3,10.10) (4,9.82)
                    (5,8.92) (6,9.58)  (7,10.28) (8,9.74)  (9,10.09)};
            \legend{Train, Validation, Test}
            % đường tham chiếu 10% (phân bố đều lý tưởng)
            \draw[dashed] (axis cs:\pgfkeysvalueof{/pgfplots/xmin},10) --
            (axis cs:\pgfkeysvalueof{/pgfplots/xmax},10);
        \end{axis}
    \end{tikzpicture}
    \caption{Phân bố lớp của các tập train, validation và test. Đường nét đứt là
        mức 10\% của phân bố đều.}
    \label{fig:class-dist}
\end{figure}

\subsection{Tiền xử lý và augmentation}
\subsection{Quy trình huấn luyện}
\subsection{Độ đo đánh giá}
\subsection{Khả năng tái lập}

\section{E1: Softmax, MLP và CNN}\label{sec:e1}
\subsection{Mục tiêu và thiết lập}
\subsection{Mô hình}
\subsection{Kết quả và đường cong học}
\subsection{Ảnh hưởng của augmentation}
\subsection{Phân tích lỗi}
\subsection{Thảo luận và kết luận ngắn}

\section{E2: Transformer (MSA và tokenization)}\label{sec:e2}
\subsection{Mục tiêu và thiết lập}
\subsection{Mô hình}
\subsection{Kiểm tra tính đúng đắn của MSA tự viết}
\subsection{Kết quả: MSA tự viết và MSA PyTorch}
\subsection{So sánh các cách tokenization}
\subsection{Đường cong học}
\subsection{Thảo luận và kết luận ngắn}

\section{E3: LSTM / GRU trên chuỗi ảnh}\label{sec:e3}
% \subsection{Mục tiêu và thiết lập}
% \subsection{Mô hình}
% \subsection{Kết quả}
% \subsection{Đường cong học}
% \subsection{Phân tích phụ thuộc thời gian}
% \subsection{Thảo luận và kết luận ngắn}

\section{So sánh tổng thể và kết luận}\label{sec:conclusion}
\subsection{Bảng so sánh tổng hợp}
\subsection{Kết luận}
\subsection{Hạn chế và hướng phát triển}

\section*{Tài liệu tham khảo}
\addcontentsline{toc}{section}{Tài liệu tham khảo}

\appendix
\section{Phụ lục}\label{sec:appendix}
\subsection{Hyperparameters đầy đủ}
\subsection{Cấu trúc repository và lệnh chạy}
\subsection{Hình bổ sung}
\subsection{Chi tiết kiểm tra MSA}
\end{document}
